\section{Conclusão Parcial}

Os resultados obtidos nesta etapa da pesquisa permitem concluir que a metodologia de monitoramento via GTFS Realtime, integrada ao tratamento estatístico por Modelos Lineares de Efeitos Mistos (LMM), é eficaz para diagnosticar a precariedade da pontualidade no transporte público de Belo Horizonte. A aplicação do modelo à linha 4110 revelou que o sistema opera sob um atraso basal crítico de aproximadamente 15,5 minutos, o qual é significativamente agravado durante o horário de pico (16h às 19h), sofrendo um acréscimo médio de 9,4 minutos.

Cientificamente, a pesquisa validou a necessidade de considerar a variância intragrupo (viagens) como um componente essencial da análise. O isolamento desses efeitos aleatórios permitiu constatar que o atraso não é um fenômeno homogêneo, mas sim uma variável complexa influenciada tanto por fatores sistêmicos (congestionamento urbano) quanto por variáveis estocásticas da operação diária.

Mesmo diante das limitações técnicas identificadas, como as lacunas na transmissão de dados GPS e instabilidades na API de dados abertos, o processo de interpolação linear espacial mostrou-se robusto o suficiente para gerar indicadores de desempenho com alto grau de confiança estatística ($p < 0,001$).

Em suma, as evidências coletadas até o momento apontam para um descompasso estrutural entre o planejamento teórico da BHTRANS e a realidade operacional da linha estudada. Esta conclusão parcial fundamenta a necessidade de uma revisão dos tempos de ciclo e sugere que a tecnologia de dados abertos, quando submetida ao rigoroso tratamento econométrico proposto neste trabalho, constitui uma ferramenta poderosa e de baixo custo para a gestão e fiscalização da mobilidade urbana em grandes metrópoles brasileiras.